\section*{ЗАКЛЮЧЕНИЕ}
\addcontentsline{toc}{section}{ЗАКЛЮЧЕНИЕ}

В ходе выполнения выпускной квалификационной работы был разработан модуль взаимодействия с пользователем для онлайн-площадки аренды вещей. Разработанный модуль включает подсистемы чатов, уведомлений, формирования договоров аренды и демонстрационного подписания договоров с использованием электронной подписи.

В процессе работы были рассмотрены особенности организации пользовательского взаимодействия в сервисах онлайн-аренды, определены основные сценарии коммуникации между арендатором и арендодателем, а также проанализированы требования к уведомлениям, истории переписки и электронному документообороту. На основе проведенного анализа были сформированы требования к программным модулям и определены основные API-маршруты серверной части.

Основные результаты работы:

\begin{enumerate}
\item Проведен анализ предметной области онлайн-площадок аренды вещей в части коммуникации пользователей, уведомлений и документооборота.
\item Разработана концептуальная модель модуля взаимодействия с пользователем, включающая подсистемы чатов, уведомлений и договоров аренды.
\item Спроектирована архитектура серверной части модуля, определены модели данных, API-маршруты и правила доступа к данным.
\item Реализована подсистема чатов, обеспечивающая создание диалогов между пользователями, хранение сообщений, отправку изображений и отметку входящих сообщений как прочитанных.
\item Реализована подсистема уведомлений, обеспечивающая создание системных уведомлений при событиях бронирования, оплаты, возврата и получения нового сообщения.
\item Реализована подсистема договоров, обеспечивающая автоматическое формирование PDF-договора аренды по данным бронирования и демонстрационное подписание договора арендатором и арендодателем.
\item Выполнено тестирование разработанных модулей средствами Django и Django REST Framework. Проверены основные пользовательские сценарии, корректность REST API, создание уведомлений, работа сообщений и подписание договоров.
\end{enumerate}

Все требования, объявленные в техническом задании, были реализованы. Разработанные модули обеспечивают серверную логику взаимодействия пользователей внутри онлайн-площадки аренды вещей и могут быть использованы как часть единой программно-информационной системы.

Готовый рабочий проект представлен Django-приложениями и набором API для работы с чатами, уведомлениями, договорами аренды и электронной подписью. Серверная часть проекта подготовлена для интеграции с пользовательским интерфейсом и дальнейшего расширения функциональности системы.